\section{\heiti 结论}
\label{sec:conclusion}

本文提出循环读出压缩：编码器在输入序列中携带 $K$ 个可学习特殊 token，以像素重建为训练信号把整幅图像压缩到这些 token 上；解码器以权重共享的 Transformer 块反复读出，第 $t$ 步只读取其读窗口，查询由固定查询基与上一步输出相加得到，每一步都直接预测整幅图像。该结构使 token 预算与读出步数绑定，同一模型可在推理时停在任意一步。结构上它受自适应长度分词（ALIT）的循环分配\cite{duggal2025alit}启发：同样以“权重共享的循环 + 每步容量递增”换取一族长度不同的码，但循环被移到解码侧，容量增长由确定性的分块窗口给出（第~\ref{sec:alit}~节）。

在 ImageNet-100 上的受控实验给出三点结论。第一，在 240M 档（编码器与解码器各约 240 M，合计 481.54 M）上，以 token 数计的压缩率由旧 63M 基准的 3.06 倍（196$\to$64）提高到 6.12 倍（196$\to$32），单步读出达到 32.35 dB，显著高于逐 patch 直流基线（18.45 dB），说明 32 个特殊 token 可以保留整幅图像的主要信息；把同一模型展开到 16 步，PSNR 由 22.07 dB 升到 27.94 dB。与逐块 PCA 的同值数对照按预算分层：在 $\le3{,}584$ 个连续值的紧预算端，本文 448 维与 224 维瓶颈码均高于同值数 PCA（$+0.89\sim+1.41$ dB）；7{,}168 值处 448 维码仍高 0.34 dB（224 维码仅 L1 更低）；预算放宽到 14{,}336 值及以上时 PCA 反超（$+1.37$ dB；主模型在 39{,}936 值处为 27.94 dB 对 35.75 dB）——逐块 PCA 保留 196 个空间位置，预算越松越占优。把每个特殊 token 由 1248 维压到 448 维（每图连续值降为 36$\%$）只损失 0.25 dB，且在同一 14{,}336 值下比旧 63M 基准高 3.54 dB，说明这类码的冗余主要在每个 token 的宽度方向。第二，人评复杂度与重建误差在三个子集上均正相关（$\rho=+0.33\sim+0.54$，5 个容量档与 3 个子集共 15 个组合全部为正），容量缩小的是复杂度惩罚的绝对幅度而非相对结构；SAVOIAS 上的分布外结果与 ImageNet-100 验证集接近，最优停步同为第 15 步。第三，容量阶梯在 240M 档达到峰值（27.94 dB），继续加宽到 720.6 M 反而回落到 27.31 dB，240M 是本协议下的操作点；每步 2 个 token、16 步的读出损失（相对单步读满 32 个 token）在 4.10--4.66 dB 之间，且几乎不随容量变化。

本文的结论限定在“重建”这一信息保持口径内。与现有分词器在同一协议下的逐档对比、把该压缩表示接到下游视觉—语言任务、把读出状态跨帧传递到视频流、引入量化与熵编码以给出真实比特率、以及在同预算下比较单步与循环，是后续工作；在参数量对齐下比较“每轮新增记忆槽”与“每步扩大读窗口”两种增容方式，则是判定本文与 ALIT 差异来源的直接实验。
